% ----------------------------------------------------------------------
%  3.6. Магнитные свойства: 3.6.1–3.6.5 (определения, атомы, диа- и
%  парамагнетизм, восприимчивость слабых магнетиков)
% ----------------------------------------------------------------------
\begin{orientir}
Магнетизм вещества~— квантовое явление: его источники~— орбитальное
движение и спин электронов. Мы начнём с определений и классификации
(3.6.1), разберём магнитные моменты атомов (3.6.2), слабый магнетизм~—
диа- и парамагнетизм (3.6.3--3.6.5), затем сильный: природу
ферромагнетизма (3.6.6), домены (3.6.7), свойства ферромагнетиков и
магнитные материалы (3.6.8), антиферро- и ферримагнетики (3.6.9).
\end{orientir}

\subsection{Основные определения. Классификация веществ}\label{ssec:magn-def}

Любое тело во внешнем магнитном поле $\vb H$ \emph{намагничивается}~—
приобретает магнитный момент. Магнитный момент единицы объёма~—
\term{намагниченность}:
\begin{equation}\label{eq:magnetization}
\vb I=\frac{\sum_i\vb p_i}{V},\qquad \vb I=\chi\vb H ,
\end{equation}
где $\vb p_i$~— магнитные моменты атомов, $\chi$~— безразмерная
\term{магнитная восприимчивость} (объёмная). В анизотропных телах
$\chi$~— тензор, и $\vb I$ может быть не параллельна $\vb H$. Используют
также удельную $\chi_\text{г}=\chi/\rho$ и атомную (молярную)
$\chi_a=\chi A/\rho$ восприимчивости. Индукция внутри тела
\begin{keyeq}\label{eq:B-mu}
\vb B=\mu_0(\vb H+\vb I)=\mu_0(1+\chi)\vb H=\mu_0\mu\vb H ,
\end{keyeq}
где $\mu_0=4\pi\cdot10^{-7}$~Гн/м, $\mu=1+\chi$~— относительная
магнитная проницаемость. В СИ $I$ и $H$ измеряют в А/м, $B$~— в теслах.
Восприимчивость может быть и положительной (намагниченность по полю),
и отрицательной (против поля).

\paragraph{Классификация}
Все вещества делятся на два больших класса (рис.~\ref{fig:magn-classes}).
\begin{itemize}
\item \term{Магнитно-неупорядоченные}~— \emph{диамагнетики} и
\emph{парамагнетики}. Элементарные моменты (если они есть) ориентированы
хаотично. Восприимчивость мала ($10^{-7}$--$10^{-5}$ по модулю) и не
зависит от поля; у парамагнетиков $\chi>0$, у диамагнетиков $\chi<0$.
В неоднородном поле диамагнетик выталкивается из области сильного
поля, парамагнетик втягивается~— на этом основаны измерения
восприимчивости. Их называют \emph{слабомагнитными}.
\item \term{Магнитно-упорядоченные}~— \emph{ферромагнетики},
\emph{антиферромагнетики}, \emph{ферримагнетики}. Локализованные в узлах
моменты упорядочены \emph{обменным взаимодействием} электронов (оно
электростатической природы). У ферромагнетиков моменты одинаковы и
параллельны, у антиферромагнетиков равные моменты антипараллельны и
компенсируются, у ферримагнетиков антипараллельные моменты различны и
не компенсируются. Ферро- и ферримагнетики \emph{спонтанно намагничены}
и даже без поля имеют макроскопическую намагниченность; их
восприимчивость достигает $10^5$ и зависит от поля. Порядок исчезает
выше критической температуры~— точки Кюри (ферро- и ферримагнетики) или
Нееля (антиферромагнетики).
\end{itemize}

\begin{figure}[ht]
\centering
\begin{tikzpicture}[x=1cm, y=1cm, sp/.style={-{Stealth[length=1.5mm]}, line width=0.9pt}]
  \foreach \sh/\lab in {0/диамагнетик, 2.9/парамагнетик, 5.8/ферромагнетик, 8.7/антиферромагнетик, 11.6/ферримагнетик}{
    \node[font=\scriptsize] at (\sh+0.9,-0.55) {\lab};
    \foreach \i in {0,...,3}{\foreach \j in {0,1,2}{\fill[atomA] (\sh+\i*0.6,\j*0.6+0.3) circle (0.09);}}}
  \foreach \i/\a in {0/20,1/150,2/260,3/80}{\foreach \j/\b in {0/0,1/110,2/220}{
    \pgfmathsetmacro{\ang}{\a+\b}
    \draw[sp, fmBlue] ($(2.9+\i*0.6,\j*0.6+0.3)+(\ang+180:0.2)$) -- ($(2.9+\i*0.6,\j*0.6+0.3)+(\ang:0.2)$);}}
  \foreach \i in {0,...,3}{\foreach \j in {0,1,2}{
    \draw[sp, fmRed] (5.8+\i*0.6,\j*0.6+0.1) -- (5.8+\i*0.6,\j*0.6+0.5);}}
  \foreach \i in {0,...,3}{\foreach \j in {0,1,2}{
    \pgfmathtruncatemacro{\pp}{mod(\i+\j,2)}
    \ifnum\pp=0 \draw[sp, fmRed] (8.7+\i*0.6,\j*0.6+0.1) -- (8.7+\i*0.6,\j*0.6+0.5);
    \else \draw[sp, fmBlue] (8.7+\i*0.6,\j*0.6+0.5) -- (8.7+\i*0.6,\j*0.6+0.1);\fi}}
  \foreach \i in {0,...,3}{\foreach \j in {0,1,2}{
    \pgfmathtruncatemacro{\pp}{mod(\i+\j,2)}
    \ifnum\pp=0 \draw[sp, fmRed] (11.6+\i*0.6,\j*0.6+0.06) -- (11.6+\i*0.6,\j*0.6+0.54);
    \else \draw[sp, fmBlue] (11.6+\i*0.6,\j*0.6+0.42) -- (11.6+\i*0.6,\j*0.6+0.18);\fi}}
\end{tikzpicture}
\caption{Расположение атомных магнитных моментов (стрелки) в отсутствие
поля: у диамагнетика собственных моментов нет; у парамагнетика они
хаотичны; у ферромагнетика параллельны; у антиферромагнетика
антипараллельны и равны; у ферримагнетика антипараллельны и не равны}
\label{fig:magn-classes}
\end{figure}

\subsection{Магнитные свойства свободных атомов}\label{ssec:magn-atoms}

\paragraph{Орбитальный момент}
Электрон, обращающийся по орбите, подобен витку с током $i$ площадью
$s$: его магнитный момент $M=is$. При частоте обращения $\omega_0$ ток
$i=e\omega_0/2\pi$, и
\[
M=\frac{e\omega_0}{2\pi}\pi r^2=\frac e{2m}\,mr^2\omega_0=-\frac e{2m}\,l_\text{мех},
\]
где $l_\text{мех}=mr^2\omega_0$~— механический (угловой) момент; знак
минус~— из-за отрицательного заряда электрона. Отношение магнитного
момента к механическому~— \term{гиромагнитное отношение}~— для
орбитального движения $\gamma=e/2m$\ispr{В оригинале (с.~524)
гиромагнитным отношением названо $l_\text{мех}/M$, но приравнено к $e/2m$.
Отношение магнитного момента к механическому $M/l_\text{мех}=e/2m$.}.
Механический момент квантуется с квантом $\hbar$, отсюда квант
магнитного момента~— \term{магнетон Бора}:
\begin{keyeq}\label{eq:bohr-magneton}
M_\text{Б}=\frac{e\hbar}{2m}=9{,}274\cdot10^{-24}~\text{А}\cdot\text{м}^2 .
\end{keyeq}

\paragraph{Спиновый момент}
Электрон имеет и собственный (\term{спиновый}) магнитный момент,
численно равный магнетону Бора, хотя механический спиновый момент
равен лишь $\hbar/2$: спиновое гиромагнитное отношение вдвое больше
орбитального. Протоны и нейтроны тоже обладают моментами, но в
$\sim1836$ раз меньшими (ядерный магнетон $5{,}051\cdot10^{-27}$~А$\cdot$м$^2$);
для магнетизма твёрдых тел ими можно пренебречь.

\paragraph{Многоэлектронный атом}
В квантовой механике модули моментов выражаются через квантовые числа:
орбитальный~— $\sqrt{l(l+1)}\,\hbar$ и $\sqrt{l(l+1)}\,M_\text{Б}$
($s$-электроны с $l=0$ орбитального момента не имеют), спиновый~—
$\sqrt{s(s+1)}\,\hbar$ и $2\sqrt{s(s+1)}\,M_\text{Б}=\sqrt3M_\text{Б}$. У
многоэлектронного атома сначала складывают орбитальные моменты в
$\vb L$ и спиновые в $\vb S$, затем~— в полный момент $\vb J=\vb L+\vb S$
(\term{$LS$-связь}, связь Рассела--Саундерса). Полный магнитный момент
атома характеризуют максимальной проекцией $gJM_\text{Б}$, где
\begin{equation}\label{eq:lande}
g=1+\frac{J(J+1)+S(S+1)-L(L+1)}{2J(J+1)}
\end{equation}
\term{множитель Ланде} ($g=2$ для чисто спинового и $g=1$ для чисто
орбитального момента).

\begin{glavnoe}
Заполненные оболочки имеют нулевой суммарный момент. По \emph{правилу
Хунда} наименьшую энергию имеет состояние с наибольшим суммарным спином
(спины параллельны) и наибольшим $L$. Поэтому у переходных металлов и
лантаноидов недостроенные $d$- и $f$-оболочки всегда обладают
нескомпенсированными моментами~— это и есть «магнитные» атомы. У
непереходных элементов магнитно активны лишь недостроенные валентные
оболочки.
\end{glavnoe}

Поле действует на атом двояко. \emph{Поляризационный} эффект~—
появление момента, направленного против поля,~— это
\term{диамагнетизм}; он присущ всем атомам. \emph{Ориентационный}
эффект~— поворот собственного момента атома по полю~— это
\term{парамагнетизм}; он есть только у атомов с ненулевым моментом.

\subsection{Физическая природа диамагнетизма}\label{ssec:diamagn}

Электронная орбита~— виток с током. По закону Ленца при внесении витка
в поле в нём индуцируется ток, магнитный момент которого направлен
против поля. Поэтому диамагнетизм есть у всех веществ, но обычно он
слабее других эффектов; в чистом виде его видно у веществ с нулевыми
собственными моментами атомов (благородные газы).

\paragraph{Ларморова прецессия}
Пусть электрон движется по круговой орбите радиуса $r$, а поле $B_0$
перпендикулярно её плоскости.
\begin{shagi}
\item Без поля центростремительную силу создаёт кулоновское
притяжение: $m\omega_0^2r=e^2/(4\pi\varepsilon_0r^2)$.
\item В поле добавляется сила Лоренца $\pm e\omega rB_0$, и
$\omega^2=\omega_0^2\pm e\omega B_0/m$.
\item Так как изменение частоты мало, $(\omega-\omega_0)(\omega+\omega_0)\approx2\omega(\omega-\omega_0)$, и
\end{shagi}
\begin{keyeq}\label{eq:larmor}
\omega=\omega_0\pm\omega_L,\qquad \omega_L=\frac{eB_0}{2m}.
\end{keyeq}
Все электроны, независимо от направления вращения, меняют частоту на
\term{ларморову частоту} $\omega_L$, и добавочный момент направлен против
поля. Если поле не перпендикулярно орбите, орбита \emph{прецессирует}
вокруг поля с той же частотой $\omega_L$ (рис.~\ref{fig:larmor}).

\begin{figure}[ht]
\centering
\begin{tikzpicture}[scale=0.9]
  \draw[-{Stealth[length=2.2mm]}, line width=1pt, fmGreen!60!black] (0,-0.2) -- (0,2.6) node[above, font=\small]{$\vb B_0$};
  \draw[fmGray, dashed] (0,1.7) ellipse (1.0 and 0.25);
  \draw[-{Stealth[length=2mm]}, line width=0.9pt, fmBlue] (0,0) -- (0.9,1.5) node[right, font=\small]{нормаль к орбите};
  \draw[fmBlue] (0.95,1.62) ellipse (0.5 and 0.18);
  \draw[-{Stealth[length=1.6mm]}, fmGray] (-0.6,1.48) arc (200:260:0.9 and 0.22);
  \node[font=\scriptsize, fmGray] at (-1.25,1.6) {$\omega_L$};
  \fill (0,0) circle (1.5pt);
\end{tikzpicture}
\caption{Ларморова прецессия: в магнитном поле нормаль к плоскости
электронной орбиты описывает конус вокруг направления поля с частотой
$\omega_L=eB_0/2m$; добавочное круговое движение создаёт момент против
поля}
\label{fig:larmor}
\end{figure}

\paragraph{Диамагнитная восприимчивость}
Добавочный момент одного электрона $\Delta M=-e^2\rho^2B_0/4m$, где
$\rho^2$~— квадрат расстояния электрона от оси, проходящей через ядро
вдоль поля. Для сферически симметричного атома
$\langle\rho^2\rangle=\langle x^2+y^2\rangle=\frac23\langle r^2\rangle$. Для
$n$ атомов в единице объёма с $z$ электронами каждый:
\begin{keyeq}\label{eq:chi-dia}
\chi_\text{диа}=-\frac{\mu_0nze^2\langle r^2\rangle}{6m}.
\end{keyeq}
Диамагнитная восприимчивость \emph{не зависит от температуры} и растёт
с числом электронов и размером орбит. Атомная диамагнитная
восприимчивость при $\langle r^2\rangle\approx10^{-20}$~м$^2$ составляет
$\chi_a\approx-4\pi\cdot2{,}8\cdot10^{-9}\,z$~м$^3$/кмоль.

\subsection{Физическая природа парамагнетизма}\label{ssec:paramagn}

\paragraph{Почему без поля нет намагниченности}
Энергия взаимодействия двух атомных моментов порядка магнетона Бора на
расстоянии $r\approx0{,}1$~нм:
\[
E_\text{маг}\approx\frac{\mu_0}{4\pi}\frac{M_\text{Б}^2}{r^3}=
10^{-7}\cdot\frac{(9{,}27\cdot10^{-24})^2}{(10^{-10})^3}\approx10^{-23}~\text{Дж},
\]
что в сто раз меньше тепловой энергии при комнатной температуре
($kT\approx4\cdot10^{-21}$~Дж). Тепловое движение разрушает ориентацию,
и без поля намагниченность парамагнетика равна нулю.

\paragraph{Теория Ланжевена--Бриллюэна}
В поле устанавливается динамическое равновесие: поле ориентирует моменты,
тепловое движение их разупорядочивает. Проекция момента на поле
квантуется: $\mu_H=gM_\text{Б}M_J$, где $M_J=-J,\ldots,+J$. Вероятность каждой
ориентации задаёт распределение Больцмана $\propto\exp(\mu_HB_0/kT)$.
Усреднение даёт
\begin{keyeq}\label{eq:brillouin}
I=ngJM_\text{Б}\,B_J(\beta),\qquad \beta=\frac{gJM_\text{Б}B_0}{kT},
\end{keyeq}
где $B_J(\beta)=\frac{2J+1}{2J}\operatorname{cth}\frac{(2J+1)\beta}{2J}-\frac1{2J}\operatorname{cth}\frac\beta{2J}$~—
\term{функция Бриллюэна}. При $J\to\infty$ (непрерывные ориентации) она
переходит в классическую \term{функцию Ланжевена}
$L(\beta)=\operatorname{cth}\beta-1/\beta$ (рис.~\ref{fig:langevin}).
\begin{itemize}
\item В сильных полях и при низких температурах ($\beta\gg1$) $B_J\to1$:
намагниченность \emph{насыщается}, $I_\text{нас}=ngJM_\text{Б}$~— все
моменты выстроены.
\item В слабых полях ($\beta\ll1$) $L(\beta)\approx\beta/3$, и
\end{itemize}
\begin{keyeq}\label{eq:curie}
\chi=\frac{n\mu_0\mu_\text{ат}^2}{3kT}=\frac CT\qquad\text{(закон Кюри)}.
\end{keyeq}

\begin{figure}[ht]
\centering
\begin{subfigure}[b]{0.48\textwidth}\centering
\begin{tikzpicture}
\begin{axis}[width=0.95\textwidth, height=4.8cm, xmin=0, xmax=6, ymin=0, ymax=1.08,
  xlabel={$\beta$}, ylabel={$I/I_\text{нас}$}, axis lines=left,
  legend style={font=\scriptsize, at={(0.97,0.03)}, anchor=south east, draw=none, fill=none}]
  \addplot[fmRed, thick, domain=0.001:6, samples=80] {tanh(x)};
  \addlegendentry{{$J=1/2$, $\mathrm{th}\,\beta$}}
  \addplot[fmBlue, thick, domain=0.001:6, samples=80] {cosh(x)/sinh(x)-1/x};
  \addlegendentry{{$J\to\infty$, $L(\beta)$}}
  \addplot[fmGray, dashed, domain=0:2.4] {x/3};
  \node[font=\scriptsize, fmGray, anchor=west] at (axis cs:2.4,0.82) {$\beta/3$};
\end{axis}
\end{tikzpicture}
\caption{}
\end{subfigure}\hfill
\begin{subfigure}[b]{0.48\textwidth}\centering
\begin{tikzpicture}
\begin{axis}[width=\textwidth, height=4.8cm, xmin=0, xmax=3, ymin=0, ymax=3,
  xlabel={$T$}, ylabel={$1/\chi$}, axis lines=left, xtick={0,0.8}, xticklabels={0,$\theta$}, ytick=\empty]
  \addplot[fmBlue, thick, domain=0:3] {x};
  \addplot[fmRed, thick, domain=0.8:3] {x-0.8};
  \node[font=\scriptsize, fmBlue, anchor=south east] at (axis cs:2.2,2.3) {Кюри: $\chi=C/T$};
  \node[font=\scriptsize, fmRed, anchor=north west] at (axis cs:1.4,0.5) {Кюри--Вейсс: $\chi=\frac C{T-\theta}$};
\end{axis}
\end{tikzpicture}
\caption{}
\end{subfigure}
\caption{\textit{а}~— намагниченность парамагнетика в зависимости от
$\beta=\mu B/kT$: при малых $\beta$ она пропорциональна полю, при больших
насыщается; \textit{б}~— законы Кюри и Кюри--Вейсса в координатах $1/\chi$~— $T$}
\label{fig:langevin}
\end{figure}

Закон Кюри хорошо выполняется для паров металлов и парамагнитных солей.
Для твёрдых парамагнитных металлов и для ферромагнетиков выше точки
Кюри он не выполняется. Вейсс предположил, что взаимодействие
атомных моментов создаёт добавочное внутреннее поле; тогда
восприимчивость подчиняется \term{закону Кюри--Вейсса}:
\begin{keyeq}\label{eq:curie-weiss}
\chi=\frac C{T-\theta},
\end{keyeq}
где $\theta$~— парамагнитная температура Кюри (у ферромагнетиков она
немного выше ферромагнитной точки Кюри $\theta_K$).

\paragraph{Парамагнетизм электронов проводимости (Паули)}
Поле ориентирует и спины свободных электронов металла. Дорфман
заметил, что диамагнитная восприимчивость металлической меди по модулю
примерно в три раза меньше, чем у газа её ионов,~— электроны
проводимости частично компенсируют диамагнетизм ионных остовов. Паули
объяснил это квантово (рис.~\ref{fig:pauli}).
\begin{shagi}
\item В поле энергия электронов со спиновым моментом по полю
уменьшается на $M_\text{Б}B_0$, против поля~— увеличивается на столько
же: две «половинки» плотности состояний сдвигаются друг относительно
друга.
\item Электроны перетекают из подзоны «против поля» в подзону «по
полю», пока уровни Ферми не сравняются. Перетекает примерно
$\frac12g(\EF)M_\text{Б}B_0$ электронов, и избыток моментов по полю равен
$g(\EF)M_\text{Б}^2B_0$.
\item В поле $B_0\sim1$~Тл $M_\text{Б}B_0\approx10^{-23}$~Дж~— много меньше
$\EF\approx10^{-19}$--$10^{-18}$~Дж: в перераспределении участвуют лишь
электроны у уровня Ферми.
\end{shagi}
\begin{keyeq}\label{eq:pauli}
\chi_\text{Паули}=\mu_0g(\EF)M_\text{Б}^2=\frac{3n\mu_0M_\text{Б}^2}{2\EF}.
\end{keyeq}
Эта восприимчивость положительна и \emph{не зависит от температуры}~— в
отличие от закона Кюри: снова работает лишь малая доля электронов у
поверхности Ферми. Кроме того, поле искривляет траектории свободных
электронов; квантование этого периодического движения даёт
\term{диамагнетизм Ландау}, равный $-\frac13$ парамагнетизма Паули. В
сумме электроны проводимости парамагнитны:
$\chi_\text{эл}=n\mu_0M_\text{Б}^2/\EF$.

\begin{figure}[ht]
\centering
\begin{tikzpicture}
\begin{axis}[width=0.55\textwidth, height=5.4cm, xmin=-1.25, xmax=1.25, ymin=0, ymax=1.5,
  xlabel={$g_-(E)$ \hspace{14mm} $g_+(E)$}, ylabel={$E$}, axis lines=middle, xtick=\empty, ytick=\empty,
  every axis x label/.style={at={(axis description cs:0.5,-0.08)}, anchor=north}]
  \addplot[fill=fmBlue!25, draw=none, domain=0:2, samples=161] ({x<1 ? sqrt(x) : 0}, {x<1 ? 0.12+x : 1.12-(x-1)});
  \addplot[fill=fmRed!25, draw=none, domain=0:2, samples=161] ({x<1 ? -sqrt(0.84*x) : 0}, {x<1 ? 0.28+0.84*x : 1.12-0.84*(x-1)});
  \addplot[fmBlue, thick, domain=0.12:1.45, samples=60] ({sqrt(x-0.12)}, {x});
  \addplot[fmRed, thick, domain=0.28:1.45, samples=60] ({-sqrt(x-0.28)}, {x});
  \draw[dashed] (axis cs:-1.25,1.12) -- (axis cs:1.25,1.12) node[above left, font=\scriptsize]{$\EF$};
  \draw[dimline] (axis cs:0.2,0.12) -- (axis cs:0.2,0.28);
  \node[font=\scriptsize, anchor=west] at (axis cs:0.22,0.2) {$2M_\text{Б}B_0$};
  \node[font=\scriptsize, fmBlue] at (axis cs:0.55,1.32) {спин по полю};
  \node[font=\scriptsize, fmRed] at (axis cs:-0.6,1.32) {против поля};
\end{axis}
\end{tikzpicture}
\caption{Парамагнетизм Паули: плотности состояний электронов со спинами
по полю ($g_+$) и против поля ($g_-$) сдвинуты в поле на $2M_\text{Б}B_0$;
при общем уровне Ферми электронов со спином по полю больше (сдвиг сильно
преувеличен)}
\label{fig:pauli}
\end{figure}

\paragraph{Адиабатическое размагничивание}
Поле упорядочивает моменты парамагнетика, тепловое движение их
разупорядочивает. На этом основан способ получения сверхнизких
температур: парамагнитную соль охлаждают жидким гелием в сильном поле,
затем изолируют и медленно размагничивают. Разупорядочение моментов
требует тепла, его отбирают у решётки, и температура падает до
$\sim10^{-3}$~К.

\subsection{Магнитная восприимчивость слабых магнетиков}\label{ssec:weak-magn}

Диамагнетизм есть у всех тел; если тело парамагнитно, значит,
парамагнитный эффект (атомов, ионов или электронов проводимости)
перекрывает диамагнитный.

\emph{Диамагнетики:}
\begin{enumerate}
\item благородные газы в любом агрегатном состоянии (заполненные
оболочки, нулевой момент);
\item газы из атомов и молекул без результирующего момента (H$_2$, N$_2$)\ispr{В
оригинале (с.~538) в качестве примера диамагнитного газа приведён
O$_2$, но молекула кислорода парамагнитна (два неспаренных электрона),
что и указано в оригинале ниже.};
\item ионные кристаллы из ионов с оболочками благородных газов (NaCl,
KI, MgCl$_2$);
\item некоторые непереходные металлы и металлоиды, у которых диамагнетизм
ионных остовов сильнее парамагнетизма электронов проводимости. Этому
способствуют большое число электронов и большие радиусы орбит
(формула \eqref{eq:chi-dia}) и малая плотность состояний на уровне Ферми
(формула \eqref{eq:pauli}). Самые сильные диамагнетики~— тяжёлые
металлы Au, Sb, Bi, Hg, Pb (табл.~\ref{tab:chi-metals}).
\end{enumerate}

\begin{table}[ht]
\centering
\caption{Объёмная магнитная восприимчивость непереходных металлов при
комнатной температуре, $\chi\cdot10^6$ (СИ)}
\label{tab:chi-metals}
\small
\begin{tabular}{@{}lcc@{\qquad}lcc@{}}
\toprule
Металл & $Z$ & $\chi\cdot10^6$ & Металл & $Z$ & $\chi\cdot10^6$ \\
\midrule
Li & 3 & $+14$ & Mg & 12 & $+11{,}9$ \\
Na & 11 & $+8{,}5$ & Zn & 30 & $-15{,}6$ \\
K & 19 & $+5{,}7$ & Cd & 48 & $-19{,}1$ \\
Cu & 29 & $-9{,}5$ & Al & 13 & $+21{,}0$ \\
Ag & 47 & $-26{,}4$ & In & 49 & $-10{,}1$ \\
Au & 79 & $-36{,}4$ & Pb & 82 & $-17{,}1$ \\
Be & 4 & $-23{,}0$ & Bi & 83 & $-163$ \\
\bottomrule
\end{tabular}
\end{table}
\noindent Значения для щелочных металлов исправлены\ispr{В табл.~3.22
оригинала для Li указано $+2{,}4$, а значения для Na и K переставлены
($+5{,}6$ и $+8{,}5$). По справочным данным объёмная восприимчивость (СИ)
Li~— около $+14\cdot10^{-6}$, Na~— $+8{,}5\cdot10^{-6}$, K~— $+5{,}7\cdot10^{-6}$.}.

\emph{Парамагнетики:}
\begin{enumerate}
\item вещества из атомов и молекул с собственным моментом: с нечётным
числом электронов (свободные атомы щелочных металлов, молекула NO,
свободные радикалы); атомы и ионы с недостроенными $d$- и
$f$-оболочками (переходные и редкоземельные элементы); некоторые
молекулы с чётным числом электронов (O$_2$, S$_2$);
\item металлы, у которых парамагнетизм электронов сильнее диамагнетизма
остовов (K, Al);
\item дефекты с нечётным числом электронов ($F$-центры в
щёлочно-галоидных кристаллах, вакансии и дивакансии в кремнии).
\end{enumerate}
